% Overleaf compiler: XeLaTeX
\documentclass[UTF8]{ctexart}

\usepackage[a4paper,margin=1in]{geometry}
\usepackage{amsmath,amssymb,bm}
\usepackage{graphicx}
\usepackage{booktabs}
\usepackage{array}
\usepackage{float}
\usepackage{hyperref}
\usepackage{enumitem}
\usepackage{xcolor}
\usepackage{iftex}

\ifPDFTeX
    \PackageError{compiler}{Please compile this file with XeLaTeX on Overleaf}{Set Menu -> Compiler -> XeLaTeX.}
\fi

\hypersetup{
    colorlinks=true,
    linkcolor=blue,
    urlcolor=blue,
    citecolor=blue
}

\newcommand{\safeincludegraphics}[2][]{%
    \IfFileExists{#2}{%
        \includegraphics[#1]{#2}%
    }{%
        \IfFileExists{mtj_ising_application_demos_output/#2}{%
            \includegraphics[#1]{mtj_ising_application_demos_output/#2}%
        }{%
            \fbox{%
                \parbox{0.86\textwidth}{%
                    \centering
                    Missing figure file:\\[0.5em]
                    \texttt{\detokenize{#2}}\\[0.5em]
                    Upload this PNG to Overleaf.
                }%
            }%
        }%
    }%
}

\title{MTJ K-State Ising Machine Application Demos}
\author{}
\date{}

\begin{document}
\maketitle

\section*{Summary}

本文总结当前基于 MTJ physical K-state stochastic model 的 Ising machine 应用 demo。核心思路是：先将实际组合优化问题写成 QUBO，再转换为 Ising 能量函数，最后用从 MTJ 实验数据拟合出的 K-state 隐状态、dwell rate、transition behavior 和随机序列来驱动随机更新。

相关脚本与输出包括：
\begin{itemize}[leftmargin=2em]
    \item 主脚本：\texttt{mtj\_ising\_application\_demos.py}
    \item 总结文件：\texttt{mtj\_ising\_application\_demos\_output/application\_demo\_summary.json}
    \item TSP scaling：\texttt{mtj\_ising\_application\_demos\_output/tsp\_scaling\_summary.csv}
    \item Max-Cut scaling：\texttt{mtj\_ising\_application\_demos\_output/maxcut\_scaling\_summary.csv}
    \item Portfolio scaling：\texttt{mtj\_ising\_application\_demos\_output/portfolio\_scaling\_summary.csv}
\end{itemize}

\section{Physical K-State Ising Solver}

当前模型先从 measured MTJ time-series 中提取多状态 stochastic profile。对 good dataset 的拟合结果为：

\begin{table}[H]
\centering
\begin{tabular}{lr}
\toprule
Quantity & Value \\
\midrule
Selected $K$ & 4 \\
p-bit dimensions & 2 \\
State entropy & 1.665 bits \\
\bottomrule
\end{tabular}
\caption{MTJ K-state profile summary.}
\end{table}

该 K-state profile 提供三个核心硬件随机资源：
\begin{enumerate}[leftmargin=2em]
    \item 离散隐状态序列：\texttt{step\_states}
    \item 每个状态的 dwell/exit rate：\texttt{exit\_rates}
    \item 由 MTJ 信号重采样得到的随机数矩阵：\texttt{random\_u}
\end{enumerate}

优化问题统一写为 Ising energy：
\begin{equation}
E(\bm{s}) =
-\frac{1}{2}\bm{s}^{T}J\bm{s}
-\bm{h}^{T}\bm{s}
+ C,
\qquad
s_i\in\{-1,+1\}.
\end{equation}

每次更新一个 spin，局部场为：
\begin{equation}
H_i = \sum_j J_{ij}s_j + h_i.
\end{equation}

翻转第 $i$ 个 spin 的能量变化为：
\begin{equation}
\Delta E_i = 2s_iH_i.
\end{equation}

MTJ K-state solver 中，proposal probability 由硬件状态 bias 和 Ising 局部场共同决定：
\begin{equation}
p(s_i^{\mathrm{new}}=+1)
=
\sigma\left[
2\left(
g_{\mathrm{data}}b_k
+
g_{\mathrm{coupling}}
\frac{H_i}{E_{\mathrm{scale}}}
\right)
\right],
\end{equation}
其中
\begin{equation}
\sigma(x)=\frac{1}{1+e^{-x}}.
\end{equation}

若 proposed spin 会降低能量，则直接接受；若升高能量，则使用由 MTJ dwell/exit rate 控制的 hazard acceptance：
\begin{equation}
p_{\mathrm{accept}}
=
1-\exp\left[
-\gamma r_k\Delta t
\exp\left(
-\frac{\Delta E}{E_{\mathrm{scale}}}
\right)
\right].
\end{equation}

这里 $r_k$ 是当前 K-state 的 exit rate，$\Delta t$ 来自 MTJ profile 的时间步长，$\gamma$ 是 hazard gain。

\section{QUBO to Ising Mapping}

所有实际问题先写成 QUBO：
\begin{equation}
Q(\bm{x})=
C+\sum_i a_i x_i
+\sum_{i<j} b_{ij}x_ix_j,
\qquad
x_i\in\{0,1\}.
\end{equation}

二进制变量和 Ising spin 的关系为：
\begin{equation}
x_i=\frac{s_i+1}{2}.
\end{equation}

代入后可得：
\begin{equation}
Q(\bm{x}(\bm{s}))
=
-\frac{1}{2}\bm{s}^{T}J\bm{s}
-\bm{h}^{T}\bm{s}
+C'.
\end{equation}

因此同一个 MTJ K-state Ising solver 可以统一处理 TSP、Max-Cut、Portfolio 等不同问题。

\section{Traveling Salesman Problem}

\subsection{Problem Meaning}

TSP 表示路线规划问题：给定 $N$ 个城市和两两距离，寻找一条经过每个城市一次并回到起点的最短闭合路径。它可代表物流配送路径规划、芯片布线或探针测试路径优化、机器人巡检路径规划，以及任务访问顺序优化。

\subsection{QUBO Formulation}

定义变量：
\begin{equation}
x_{i,t} =
\begin{cases}
1, & \text{city } i \text{ is visited at position } t,\\
0, & \text{otherwise}.
\end{cases}
\end{equation}

每个位置只能有一个城市：
\begin{equation}
\sum_i x_{i,t}=1,\qquad \forall t.
\end{equation}

每个城市只能出现一次：
\begin{equation}
\sum_t x_{i,t}=1,\qquad \forall i.
\end{equation}

路径长度目标为：
\begin{equation}
L(\bm{x})
=
\sum_{t=0}^{N-1}
\sum_{i,j}
d_{ij}x_{i,t}x_{j,t+1},
\end{equation}
其中 $t+1$ 使用周期边界，即最后一个位置回到第一个位置。

完整 QUBO 为：
\begin{equation}
Q_{\mathrm{TSP}}(\bm{x})
=
L(\bm{x})
+
A\sum_t
\left(
\sum_i x_{i,t}-1
\right)^2
+
A\sum_i
\left(
\sum_t x_{i,t}-1
\right)^2.
\end{equation}

为了去掉旋转等价性，demo 中固定 city 0 在 position 0：
\begin{equation}
A(x_{0,0}-1)^2.
\end{equation}

\subsection{Results}

TSP 从 8 城市扩展到 24 城市。变量数随城市数平方增长：
\begin{equation}
n_{\mathrm{var}}=N^2.
\end{equation}

\begin{table}[H]
\centering
\resizebox{\textwidth}{!}{
\begin{tabular}{rrrrrrrrr}
\toprule
Cities & Variables & Ref. length & Best raw/ref & Median raw/ref & Best 2-opt/ref & Median 2-opt/ref & Feasible rate & Median violation \\
\midrule
8  & 64  & 23.346 & 1.255 & 1.468 & 1.000 & 1.156 & 0.458 & 1.0 \\
12 & 144 & 34.645 & 1.520 & 1.762 & 1.000 & 1.091 & 0.000 & 10.5 \\
16 & 256 & 48.033 & 1.760 & 2.335 & 1.000 & 1.120 & 0.000 & 19.0 \\
20 & 400 & 59.948 & 2.044 & 2.584 & 1.000 & 1.170 & 0.000 & 40.5 \\
24 & 576 & 71.952 & 1.838 & 2.878 & 1.026 & 1.260 & 0.000 & 57.5 \\
\bottomrule
\end{tabular}}
\caption{TSP scaling results. For 8 cities, the reference is exact optimum. For 12/16/20/24 cities, the reference is generated by nearest-neighbor all-start plus 2-opt.}
\end{table}

主要观察如下：
\begin{itemize}[leftmargin=2em]
    \item TSP 的 QUBO encoding 随规模快速变硬，因为 one-hot constraints 数量和变量数都按 $N^2$ 增长。
    \item Raw Ising decoded route 的质量随城市数增加明显下降。
    \item Raw one-hot feasible rate 在 12 城市之后降为 0，说明单 bit flip dynamics 很难直接满足 TSP 的 permutation matrix constraint。
    \item Raw Ising state 仍能作为 route seed。经过 2-opt repair 后，8/12/16/20 城市均达到 reference，24 城市达到 reference 的 $1.026\times$。
\end{itemize}

\begin{figure}[H]
\centering
\safeincludegraphics[width=0.95\textwidth]{tsp_route.png}
\caption{24-city TSP route from MTJ K-state Ising demo. Gray dashed line is raw Ising decoded route; orange line is 2-opt refined route.}
\end{figure}

\begin{figure}[H]
\centering
\safeincludegraphics[width=0.95\textwidth]{tsp_decoded_route_trend.png}
\caption{TSP decoded route quality trend from 8 to 24 cities.}
\end{figure}

\begin{figure}[H]
\centering
\safeincludegraphics[width=0.95\textwidth]{convergence.png}
\caption{TSP scaling convergence under MTJ K-state stochastic updates.}
\end{figure}

\section{Weighted Max-Cut}

\subsection{Problem Meaning}

Max-Cut 是图划分问题：给定带权图 $G=(V,E)$，将节点分成两组，使跨越两组的边权重总和最大。它可代表社区划分、网络分割、VLSI partitioning、图聚类，以及资源冲突图的二分优化。

\subsection{QUBO / Ising Formulation}

令
\begin{equation}
x_i =
\begin{cases}
0, & \text{node } i \text{ in partition A},\\
1, & \text{node } i \text{ in partition B}.
\end{cases}
\end{equation}

边 $(i,j)$ 被 cut 当且仅当 $x_i\ne x_j$。对二进制变量，有：
\begin{equation}
\mathrm{cut}_{ij}
=
x_i+x_j-2x_ix_j.
\end{equation}

最大化 cut weight：
\begin{equation}
\max_{\bm{x}}
\sum_{(i,j)\in E}
w_{ij}
\left(
x_i+x_j-2x_ix_j
\right).
\end{equation}

转换成 QUBO minimization：
\begin{equation}
Q_{\mathrm{MaxCut}}(\bm{x})
=
-\sum_{(i,j)\in E}
w_{ij}
\left(
x_i+x_j-2x_ix_j
\right).
\end{equation}

Max-Cut 没有 one-hot 约束，因此比 TSP 更适合直接由 Ising dynamics 表示。

\subsection{Results}

\begin{table}[H]
\centering
\begin{tabular}{rrrrrrr}
\toprule
Nodes & Variables & Edges & Ref. cut & Best ratio & Median ratio & Hit rate \\
\midrule
10 & 10 & 16  & 47  & 1.000 & 1.000 & 1.000 \\
16 & 16 & 53  & 154 & 1.000 & 1.000 & 0.833 \\
24 & 24 & 108 & 313 & 1.000 & 1.000 & 0.542 \\
32 & 32 & 130 & 380 & 1.000 & 0.987 & 0.167 \\
\bottomrule
\end{tabular}
\caption{Max-Cut scaling results. 10/16-node references are exact; 24/32-node references are multi-start 1-flip local-search references.}
\end{table}

主要观察如下：
\begin{itemize}[leftmargin=2em]
    \item Max-Cut 的变量数只随节点数线性增长：
    \begin{equation}
    n_{\mathrm{var}}=|V|.
    \end{equation}
    \item 没有 hard one-hot constraint，因此 raw Ising state 本身就是合法 partition。
    \item 在 10/16/24 nodes 上，median result 已接近或达到 reference。
    \item 32 nodes 时 hit rate 降到 0.167，说明规模增加后 stochastic search 的成功率开始下降，但 best run 仍能达到 reference。
    \item Current-state exploration 曲线显示优化器不是静态停住，而是在较优区域附近继续跳动。
\end{itemize}

\begin{figure}[H]
\centering
\safeincludegraphics[width=0.95\textwidth]{maxcut_convergence.png}
\caption{Structured Max-Cut convergence: best-so-far convergence, current-state exploration, and reference hit rate.}
\end{figure}

\begin{figure}[H]
\centering
\safeincludegraphics[width=0.75\textwidth]{maxcut_partition.png}
\caption{Example weighted Max-Cut partition.}
\end{figure}

\section{Cardinality-Constrained Portfolio Selection}

\subsection{Problem Meaning}

Portfolio selection 是投资组合选择问题：给定资产收益和风险，选择固定数量的资产，使收益高、风险低。它可代表金融投资组合优化、多候选设计方案选择、传感器或特征选择，以及固定预算下的资源配置。

\subsection{QUBO Formulation}

令
\begin{equation}
x_i =
\begin{cases}
1, & \text{asset } i \text{ is selected},\\
0, & \text{otherwise}.
\end{cases}
\end{equation}

收益为：
\begin{equation}
R(\bm{x})=\bm{\mu}^{T}\bm{x}.
\end{equation}

风险用 covariance matrix 表示：
\begin{equation}
\mathrm{Risk}(\bm{x})
=
\bm{x}^{T}\Sigma\bm{x}.
\end{equation}

选择固定数量 $K$ 个资产：
\begin{equation}
\sum_i x_i=K.
\end{equation}

QUBO minimization 写为：
\begin{equation}
Q_{\mathrm{Portfolio}}(\bm{x})
=
\lambda_{\mathrm{risk}}
\bm{x}^{T}\Sigma\bm{x}
-
\lambda_{\mathrm{return}}
\bm{\mu}^{T}\bm{x}
+
A
\left(
\sum_i x_i-K
\right)^2.
\end{equation}

其中 $\lambda_{\mathrm{risk}}$ 控制风险权重，$\lambda_{\mathrm{return}}$ 控制收益权重，$A$ 是 cardinality penalty。

\subsection{Results}

\begin{table}[H]
\centering
\begin{tabular}{rrrrrrr}
\toprule
Assets & Variables & Choose $K$ & Ref. objective & Best gap & Median gap & Hit rate \\
\midrule
8  & 8  & 3 & -0.506 & 0.000 & 0.004 & 0.417 \\
12 & 12 & 3 & -0.906 & 0.000 & 0.062 & 0.083 \\
16 & 16 & 4 & -0.983 & 0.025 & 0.104 & 0.000 \\
24 & 24 & 6 & -1.840 & 0.157 & 0.278 & 0.000 \\
\bottomrule
\end{tabular}
\caption{Cardinality-constrained portfolio scaling results. Feasible cardinality rate is 1.0 for all listed sizes.}
\end{table}

主要观察如下：
\begin{itemize}[leftmargin=2em]
    \item Portfolio 的变量数随资产数线性增长：
    \begin{equation}
    n_{\mathrm{var}}=N_{\mathrm{assets}}.
    \end{equation}
    \item Cardinality constraint 比 TSP 的 permutation constraint 简单，因此 feasible rate 始终为 1.0。
    \item 随资产数增加，达到 exact reference 的 hit rate 下降：8 assets 为 0.417，12 assets 为 0.083，16/24 assets 为 0。
    \item 虽然 16/24 assets 没有命中 exact reference，但 best gap 仍有限，说明 MTJ K-state solver 能找到接近最优的组合。
    \item Current-state exploration 曲线显示当前状态在低 gap 区域持续波动，说明系统仍保持随机探索能力。
\end{itemize}

\begin{figure}[H]
\centering
\safeincludegraphics[width=0.95\textwidth]{portfolio_convergence.png}
\caption{Structured portfolio convergence: best-so-far convergence, current-state exploration, and reference hit rate.}
\end{figure}

\begin{figure}[H]
\centering
\safeincludegraphics[width=0.95\textwidth]{portfolio_selection.png}
\caption{Example portfolio selection result.}
\end{figure}

\section{Cross-Problem Comparison}

三个问题展示了 Ising machine 在不同结构上的行为差异。

\begin{table}[H]
\centering
\resizebox{\textwidth}{!}{
\begin{tabular}{lllll}
\toprule
Problem & Variable growth & Constraint type & Raw Ising validity & Main bottleneck \\
\midrule
TSP & $N^2$ & permutation / one-hot & often invalid for large $N$ & hard one-hot constraints \\
Max-Cut & $|V|$ & unconstrained binary partition & always valid & larger graph search space \\
Portfolio & $N$ & cardinality & valid in current runs & objective landscape and exact hit rate \\
\bottomrule
\end{tabular}}
\caption{Cross-problem comparison.}
\end{table}

从结果看：
\begin{enumerate}[leftmargin=2em]
    \item Max-Cut 最自然地适配 Ising machine，因为问题本身就是二值 spin partition。
    \item Portfolio 也适合 QUBO/Ising，但固定 cardinality penalty 会影响大规模 exact hit rate。
    \item TSP 最难，因为它需要 permutation matrix。直接 Ising dynamics 很难维持 one-hot constraints，因此需要 decoding/repair，例如 2-opt。
\end{enumerate}

\section{Interpretation of Convergence Figures}

当前 convergence figures 的结构如下：
\begin{itemize}[leftmargin=2em]
    \item solid line：多次 runs 的 median trajectory；
    \item translucent band：25\%--75\% quantile interval；
    \item best-so-far convergence：历史最优能量随 annealing progress 的下降；
    \item current-state exploration：当前状态能量的中位数，用于观察系统是否仍在探索；
    \item hit rate：截至当前 progress，有多少 runs 达到 reference solution。
\end{itemize}

对于 TSP，convergence 图显示的是不同 city count 的 normalized QUBO gap：
\begin{equation}
g(t)
=
\frac{
E_{\mathrm{best}}(t)-E_{\mathrm{best,all}}
}{
E_{\mathrm{best}}(0)-E_{\mathrm{best,all}}
}.
\end{equation}

对于 Max-Cut 和 Portfolio，structured convergence 使用 reference objective：
\begin{equation}
g(t)
=
\frac{
E_{\mathrm{best}}(t)-E_{\mathrm{ref}}
}{
E_{\mathrm{best}}(0)-E_{\mathrm{ref}}
}.
\end{equation}

Hit rate 定义为：
\begin{equation}
\mathrm{HitRate}(t)
=
\frac{1}{M}
\sum_{m=1}^{M}
\mathbf{1}
\left[
E_{\mathrm{best}}^{(m)}(t)
-E_{\mathrm{ref}}
\le
\epsilon
\right].
\end{equation}

\section{Main Conclusions}

\begin{enumerate}[leftmargin=2em]
    \item The MTJ physical K-state model can be used as a hardware-derived stochastic driver for Ising/QUBO optimization.
    \item Problem structure strongly determines performance. Unconstrained binary problems such as Max-Cut are most compatible with direct Ising dynamics.
    \item TSP exposes the limitation of direct single-spin updates under strong one-hot constraints. The raw Ising state becomes infeasible as problem size increases, but still contains route information useful for local repair.
    \item Portfolio shows that cardinality-constrained selection remains feasible, but exact hit rate decreases with asset count.
    \item The most practical pipeline is not pure Ising decoding alone, but:
\end{enumerate}

\begin{equation}
\text{problem}
\rightarrow
\text{QUBO}
\rightarrow
\text{Ising solver}
\rightarrow
\text{domain-specific decoding/repair}.
\end{equation}

For TSP, the repair is 2-opt. For Max-Cut, no repair is needed. For Portfolio, the penalty formulation already produced feasible cardinality in these runs.

\end{document}
